% Porteurs : Nicole (comparaison des familles), Steve (métriques), Lynn (locuteur, WER,
% analyse d'erreurs). Règle E8 : chaque chiffre vient d'un fichier de training/results/, cité
% en commentaire à côté de la table.
\section{Expériences et résultats}
\label{sec:resultats}

\subsection{Métriques}
\label{sec:metriques}

\aRediger{Steve}{F1 macro (métrique principale, P4) et pourquoi plutôt que l'accuracy ;
  matrice de confusion ; WER ; EER ; intervalles de confiance par bootstrap et test de McNemar
  (E5) ; latence et RTF (E6).}

\subsection{Émotions}
\label{sec:res-emotions}

\aRediger{Nicole}{F1 macro en domaine (test de 1\,324 extraits) pour la famille A (CNN, CRNN)
  et la famille B, avec intervalles de confiance : \texttt{training/results/emotions\_family\_a\_cnn.json},
  \texttt{emotions\_family\_a\_crnn.json}, \texttt{c11\_family\_b\_emotions.json},
  \texttt{emotions\_comparisons.json}. Hors domaine (CaFE, 792 extraits) : \textbf{pas encore
  évalué}. Verdict sur H1 et H2.}

\subsection{Événements sonores (ESC-50)}
\label{sec:res-esc50}

\aRediger{Nicole}{Référence MFCC + SVM (F1 macro $0{,}403 \pm 0{,}040$ sur les 5 folds,
  \texttt{esc50\_mfcc\_svm.json}), puis familles A et B (\texttt{esc50\_family\_a\_*.json},
  \texttt{c11\_family\_b\_esc50.json}, \texttt{esc50\_comparisons.json}). Matrices en 5
  catégories dans \texttt{figures/}. Verdict sur H3 et H4.}

\subsection{Identification du locuteur}
\label{sec:res-locuteur}

% Chiffres : training/results/c15_speaker_openslr57.json (C15, GPU T4 de Kaggle).
Deux encodeurs gelés sont comparés sur le groupe d'évaluation d'OpenSLR~57 (112 locuteurs,
8\,016 phrases de test, section~\ref{sec:p11}) : ECAPA-TDNN, entraîné pour la vérification du
locuteur sur VoxCeleb, et mHuBERT-147, entraîné pour représenter la parole. La voix enregistrée
d'un locuteur est la moyenne de ses cinq vecteurs d'enregistrement, et le score d'une phrase
est sa similarité cosinus avec chaque voix enregistrée. Le seuil « inconnu » est celui de
l'EER sur le groupe de calibration ; pour mHuBERT, la couche gardée est aussi choisie sur ce
groupe.

\begin{table}[ht]
  \centering
  \small
  \begin{tabular}{@{}lrrrr@{}}
    \toprule
    Encodeur & Accuracy & EER & Connus acceptés et justes & Inconnus rejetés \\
    \midrule
    ECAPA-TDNN          & \textbf{97,7\,\%} & \textbf{2,7\,\%} & 96,8\,\% & 27,7\,\% \\
    mHuBERT-147 (couche 0) & 71,1\,\% & 10,2\,\% & 69,7\,\% & 11,0\,\% \\
    \bottomrule
  \end{tabular}
  \caption{Identification du locuteur sur OpenSLR~57 (groupe d'évaluation). Accuracy en
    ensemble fermé ; « inconnus » : phrases des locuteurs de calibration, absents des voix
    enregistrées de l'évaluation.}
  \label{tab:locuteur}
\end{table}

ECAPA-TDNN reconnaît la bonne personne dans 97,7\,\% des cas, avec peu d'écart entre
sous-corpus (\texttt{ca16} 96,2\,\%, Niger 94,4\,\%, Yaoundé 99,4\,\%). mHuBERT est nettement
moins bon, et son EER de calibration augmente de couche en couche (9,4\,\% à la couche 0,
20,8\,\% à la couche 12) : ses couches profondes gardent ce qui est dit et oublient qui parle.

En revanche, le rejet des voix inconnues est faible (27,7\,\% pour ECAPA). Le seuil est choisi
à l'EER sur des paires (phrase, locuteur) prises une à une, alors que la décision compare au
seuil le \emph{meilleur} des 112 scores, presque toujours plus élevé. Calibrer le seuil
directement sur la décision en ensemble ouvert est une perspective
(section~\ref{sec:conclusion}).

\subsection{Transcription}
\label{sec:res-wer}

\aRediger{Lynn, après C17}{WER et CER de Whisper base et small (français forcé) sur le test
  officiel d'OpenSLR~57 et sur Yaoundé, avec intervalles de confiance et RTF.}

\subsection{Analyse des erreurs}
\label{sec:res-erreurs}

\aRediger{Lynn}{Robustesse au bruit (F1 selon le SNR, E3) ; biais par genre sur le test
  équilibré des émotions (E4) ; exemples d'erreurs écoutés. \textbf{À produire.}}

\subsection{Explicabilité et efficacité}
\label{sec:res-explicabilite}

\aRediger{Nicole}{Carte Grad-CAM et test de masquage (E7) ; calibration par température et
  diagramme de fiabilité (E2) ; latence p50 et p95 et RTF (\texttt{latency\_cpu.json}).}
